\section{カラオケにおける自己開示の深さ}\label{カラオケにおける自己開示の深さ}
本節では，カラオケにおける自己開示の深さについて整理する． 

\subsection{既存の自己開示尺度}
これまでに開発された自己開示を測定できる尺度としては，榎本の「自己開示質問紙」\cite{榎本_1997_自己開示の心理学的研究}や，丹羽・丸野の「自己開示の深さを測定する尺度」\cite{自己開示尺度}が挙げられる．しかし，これらの心理尺度は，自己開示の項目や話題内容ごとに深さが定義されているのみで，音楽嗜好の共有について自己開示の深さを定義しているものではない．また，従来の研究では「音楽の好み」を1つの自己開示の項目として扱うのみで，その内部構造や深さについては既存研究では行われていない．

カラオケにおける自己開示は，曲に対する思い入れの深さ，その曲を開示するときに相手が知っているかどうかといった認知度，上手に歌えるかといった不安感などの様々な要素が影響し，一元的な尺度として示せるものではない．次項ではその構成要素について整理する．



\subsection{自己開示を構成する3軸}\label{3軸}
カラオケにおける音楽嗜好の共有を通じた自己開示を，\ref{自己開示の深さ}項で述べた榎本の整理した自己開示の深さと，伊藤\cite{カラオケにおける自動楽曲推薦}による「あなたはカラオケで選曲をする際，どんなことを考慮して選びますか？」というアンケート調査の結果より，以下の3つの軸で整理を行った．

\subsection*{1. 思い入れ}
思い入れは楽曲に対する個人的な愛着や経験の深さを表す主観的指標である．榎本が指摘するように，内面的な感情や動機の開示は自己開示の深さを特徴づける重要な要素である．思い入れの強い楽曲の選択は，その人固有の価値観や感情の表出を伴うため，より深い自己開示となる．

\subsection*{2. 歌うことの自信}
歌うことの自信は楽曲を歌唱する際の技術的な自信度を表す主観的指標である．アンケート調査において「歌いやすい，または自分の得意な曲」は，重要な選曲基準として挙げられたように，音域やリズムの難易度に対する自己評価の高いものが選ばれやすい．特に自身の歌唱能力における弱点の開示は，榎本の指摘する「弱点や欠点の開示」という深い自己開示の要素と合致する．

\subsection*{3. 人気度}
人気は楽曲の社会的認知度を表す客観的指標である．アンケート調査において「みんなが知っているかどうか」が選曲基準であるように，楽曲の認知の度合いは自己開示のしやすさを決めるうえで重要な要素である．例えば，非常にマイナーな曲を開示することは，相手に否定されるかもしれない，知らなかったら楽しめないかもしれないという否定的側面を伴うため，自己開示の深さを示す指標としてふさわしいと考えられる．

なお，テンポや歌詞の特徴といった要素は，個人の好みの差が大きく一律の指標として扱うことが困難であるため，本研究では除外する．


\subsection{自己開示を促す手法} % 段階的に自己開示するのがええんやで

% ① 玉ねぎの内側にたどり着くためには，外側から攻める必要がある
社会的浸透理論をレビューしているCarpenter\cite{carpenter_2015_SPT_Review}は，「玉ねぎモデルは，人々が対人的な相互作用を通じて他者の個人的情報の層を剥がしていき，核心に到達するというプロセスとして社会的浸透を説明する有用な比喩」と述べている．すなわち，個人の内面に関する思い入れや経験を引き出すには，その外側である一般的な音楽の好みを共有することから始める必要がある．
% ②いきなり深い自己開示をしても受容できない
\ref{受容の課題}項で述べたように，深い自己開示をいきなり行った場合，相手はそれを十分に受容することができない．カラオケにおいても，思い入れの強い曲や知名度の低い曲をいきなり選曲することは，相手の受容を得られにくい．そのため，多くの人は相手が知っている曲を選曲するなど，表面的な好みの共有に留まってしまう傾向がある．

そのため，カラオケにおける音楽嗜好の共有を通じた自己開示は，好きなアーティストや，普段歌う曲といった一般的な嗜好の共有から始まり，特定の楽曲への思い入れ，そして人生において重要な意味を持つ楽曲の共有へと，徐々に深めていくことが重要であると考えられる．このように音楽を介することで，自然な形での深い自己開示が可能となる．

しかし，このような段階的な自己開示を行う場合には，何らかの方法で自己開示を促す必要がある．次節では，カラオケ環境において段階的な自己開示を促進するための方法について検討する．